\section{系统辨识：从输入输出数据估计动力学}
\label{sec:16-Control-and-Reinforcement-Learning/01-Optimal-Control/06}

\begin{quote}
\textbf{先抓住这一点：}模型能否辨认首先由实验输入是否激励关键方向决定。一步预测拟合良好，也不保证滚动模拟或闭环控制可靠。
\end{quote}

\subsection{``模型已知''本身需要来源}

调节家中暖气前，我们未必知道房间升温多快、墙体蓄热多久或室外温度怎样渗入。控制章节写下 \(x_{t+1}=f(x_t,u_t)\) 时，常把这些参数当作已知；系统辨识追问它们怎样由输入输出数据获得。

最简单的线性离散模型是

\[
x_{t+1}=Ax_t+Bu_t+w_t.
\]

若状态都能观测，把多段数据排列为

\[
X_+=
\begin{pmatrix}A&B\end{pmatrix}
\begin{pmatrix}X_-\\U_-\end{pmatrix}
+W,
\]

便可用最小二乘估计 \(A,B\)。这一步计算熟悉，可信性却取决于实验怎样激发系统。

\subsection{没有持续激励，就无法区分参数}

若暖气输入一直保持常数，数据可能只覆盖一个狭窄温度区间，多组 \(A,B\) 都能解释观察。设计矩阵必须在相关方向上拥有足够秩和良好条件数，这通常概括为 persistent excitation。

激励不是越强越好。实体系统受舒适度、安全、磨损和能耗约束；可以使用小幅多频信号、分阶段实验或历史自然变化，但必须报告哪些动态方向从未被充分激发。参数估计区间很窄，也不能挽救结构上不可辨认的模型。

\subsection{闭环数据会把控制动作与噪声绑在一起}

在反馈控制下，\(u_t\) 由过去读数决定，而读数含噪声与扰动，因此回归变量可能与误差相关。直接最小二乘会把控制器对扰动的反应误认成被控对象的动力学。instrumental variables、prediction-error methods 或显式 closed-loop likelihood 用不同方式处理这种依赖。

这与观察数据中的混杂十分相似：动作通常根据系统状态选择，并非随机外生输入。若只记录``开了多少暖气、后来多暖''，不记录控制策略和可见历史，就难以分清动作效果与动作被触发的原因。

\subsection{状态空间模型还要同时估计隐藏状态}

只有输出

\[
y_t=Cx_t+v_t
\]

可见时，\(x_t\) 也是未知量。subspace identification 利用块 Hankel 矩阵和低秩结构恢复一个等价状态表示；EM 可在 E 步平滑 latent states，在 M 步更新参数。状态坐标本身只确定到相似变换：若 \(z_t=S^{-1}x_t\)，相应变换后的 \(A,B,C\) 产生同样输入输出分布。

因此应优先解释可预测输入输出行为、稳定模态和时间尺度，而不是给任意状态坐标赋予现实名称。

\subsection{辨识误差必须进入控制验证}

训练残差小只说明模型拟合了已见轨迹。应在不同输入段上做 one-step prediction 与 free rollout，检查创新是否近似白噪声，并观察参数随数据窗口和模型阶数是否稳定。控制器还应对模型集合或扰动范围做 robustness analysis。

系统辨识位于监督学习和控制之间：它使用回归、likelihood、Kalman smoothing 与低秩方法，却评价一个模型能否支持闭环决策。高预测精度不必意味着控制相关动态准确；反过来，一个简化模型若保留关键慢模态与增益，可能比复杂黑箱更适合安全设计。
